% Proposed B.133. Isolated owned source; central integration is root-owned.
\subsection*{The actual six-source subgroup before the cross filling}

\statement{B.133}{Proved}{The actual pre-cross six-source subgroup has its full profinite topology, and the final source kernels are specified ambient normal-closure intersections}
\label{stmt:B133}
Use the actual B.112--B.125 marking
\[
 K=F(a,B,d,e,f),\qquad c=aCa^{-1},\qquad
 Q=K\rtimes_\eta\langle j\rangle,\qquad j^2=1,
\]
and the proper pre-cross amalgam and its actual cross quotient
\begin{equation}\label{eq:B133-pre-cross}
 A=R*_{K}Q,\qquad
 r_\times=[jcj,k],\qquad
 X=A/\langle\!\langle r_\times\rangle\!\rangle_A.
\end{equation}
Here $R$ is the entire regional orbifold group, $R^0$ is its
original ordinary subgroup, $\eta$ is the actual based neck
reflection, and $k$ is the actual regional base loop.
Let $S=\langle s_1,\ldots,s_5\rangle\leq R^0$ be the actual
regional source group of B.112, with its entire selected-gap
conjugations. Put
\begin{equation}\label{eq:B133-source}
 \mathscr F=\langle js_1j,\ldots,js_5j,k\rangle\leq A.
\end{equation}
Then $\mathscr F=jSj*\langle k\rangle$ is free of rank six,
and has its full induced profinite topology in $A$. In particular
\begin{equation}\label{eq:B133-free-completion}
 \widehat{\mathscr F}\ \lhook\joinrel\longrightarrow\ \widehat A,
 \qquad
 \overline{\mathscr F}^{\,\widehat A}
   =j\widehat Sj\amalg\overline{\langle k\rangle}
   \cong\widehat F_6,
\end{equation}
where the factors are the actual closures of the specified source
words and base loop.

Let $\Gamma_0$ be the actual ordinary one-neck subgroup of $X$,
the joint kernel of the original and mirror characters. All six
words in \eqref{eq:B133-source} map into $\Gamma_0$. Define
\[
 N_\times=\langle\!\langle r_\times\rangle\!\rangle_A,
 \qquad
 \mathcal N_\times=
 \overline{\langle\!\langle r_\times\rangle\!\rangle}^{\,\widehat A}.
\]
With the embedding in \eqref{eq:B133-free-completion}, the exact
discrete and compact source kernels are
\begin{equation}\label{eq:B133-kernel-intersections}
 \begin{aligned}
 \ker(\mathscr F\longrightarrow\Gamma_0)
    &=\mathscr F\cap N_\times,\\
 \ker(\widehat{\mathscr F}\longrightarrow\widehat\Gamma_0)
    &=\widehat{\mathscr F}\cap\mathcal N_\times.
 \end{aligned}
\end{equation}
These are intersections with the actual ambient normal closures.
No equality between the closure of the first kernel and the second
kernel is asserted. The theorem does not assert injectivity of the
final six-source map, whose nonzero kernel was exhibited in B.130,
or identify its full internal normal-generating family.

\begin{proof}
\emph{B.133.1. The mirror-only cover has full regional vertices.}
Let $\mu:A\to C_2$ kill $R$ and send $j$ to the nonidentity
element. This is the mirror character, distinct from the original
ordinary character $\varepsilon$, which is nontrivial on $C$ and
trivial on $Q$. The kernel of $\mu$ is
\begin{equation}\label{eq:B133-mirror-cover}
 D=R*_{K}R',\qquad R'=jRj,
 \qquad juj=\eta(u)\quad(u\in K).
\end{equation}
To see the exact presentation, use the two cosets $1,j$ in the
ordinary Bass--Serre tree of \eqref{eq:B133-pre-cross}. They give
two $R$ vertices and one $Q\cap D=K$ vertex. Each incident edge
maps isomorphically to that middle vertex group, so collapsing
one edge gives \eqref{eq:B133-mirror-cover} with its stated
identification. The converse presentation exchanges the two
copies of $R$ by $j$ and imposes $juj=\eta(u)$ on $K$;
adjoining this involution recovers exactly $A$.
Thus $D$ has index two in $A$. Its vertices are the full groups
$R,R'$, not $R^0$ and its conjugate. Restricting the latter vertices
to the original ordinary cover produces a different graph with an
additional HNN stable letter; it is unnecessary here.

The actual source group $S$ lies in $C_R(c)$ by B.112.
The actual regional cone/neck intersection is
\begin{equation}\label{eq:B133-discrete-avoidance}
 C_R(c)\cap K=1.
\end{equation}
The base-winding character $R\to\mathbb Z$ takes $k$ to one and
all of $K$ to zero. Hence $\langle k\rangle\cap K=1$.
Conjugating \eqref{eq:B133-discrete-avoidance} by $j$, with
$jKj=K$, gives $jSj\cap K=1$. Every alternating nontrivial word
in $jSj$ and $\langle k\rangle$ is therefore reduced in the
proper amalgam \eqref{eq:B133-mirror-cover}. Since the actual
$S$ is free of rank five, this proves the discrete statement
\begin{equation}\label{eq:B133-discrete-free-source}
 \mathscr F=jSj*\langle k\rangle\cong F_6.
\end{equation}

\emph{B.133.2. The completed mirror cover is a proper amalgam.}
B.125 supplies the full induced topology of $R^0$ in $A^0$,
where $A^0=\ker\varepsilon$ has index two in $A$.
Taking finite-index cores shows that $R^0$ has its full topology
in $A$. It follows that $R$ itself has its full topology in $A$.
Indeed, for any finite-index normal subgroup $N\triangleleft R$,
choose an ambient finite-index normal subgroup $U\triangleleft A$
with $U\cap R^0\subseteq N\cap R^0$, and replace $U$ by
$U\cap A^0$. Its intersection with $R$ then lies in $N$.
The same conclusion holds for $R'=jRj$ by conjugation.

Since $D$ has finite index in $A$, its full completion embeds
in $\widehat A$. The two completions $\widehat R$ and
$\widehat R'$ embed in $\widehat D$, because their composites
with this embedding are the just established embeddings in
$\widehat A$. The actual $K$ has its full topology in $R$ by
the B.112/B.114 retraction through $R^0$ and finite index.
Finally the finite-target universal property of
\eqref{eq:B133-mirror-cover} identifies
\begin{equation}\label{eq:B133-proper-completed-cover}
 \widehat D=\widehat R\amalg_{\widehat K}\widehat R'
\end{equation}
as a proper free profinite amalgam. No residual-finiteness
assertion about an unspecified amalgam is substituted for this
actual vertex-topology argument.

\emph{B.133.3. Both actual closed source factors avoid the edge.}
We recall explicitly why the cone/neck comparison remains valid
in the full completion. The actual ordinary retraction
$r:R^0\to K$ fixes $K$ and satisfies
\begin{equation}\label{eq:B133-cone-companion}
 r(CuC)=\zeta(u),\qquad
 \zeta(a)=\zeta(B)=\zeta(d)=\zeta(e)=1,
 \quad\zeta(f)=f.
\end{equation}
Suppose $x\in\widehat K$ centralizes $C$ in $\widehat R$.
The continuous extension of \eqref{eq:B133-cone-companion}
gives $x=\widehat\zeta(x)$, so $x=f^\lambda$ for a unique
$\lambda\in\widehat{\mathbb Z}$. Uniqueness follows from the
primitive $f$-exponent retraction of the free group $K$.

The actual coarse quotient of B.112 is
\[
 R\longrightarrow
   \bigl(C_2(C)*\mathbb Z(k)\bigr)\times\mathbb Z_z,
 \qquad f\longmapsto(Ck^{-1})z.
\]
For every $n\geq2$, kill the central $z$ factor, send $C$ to
the lamp at zero, and send $k$ to the positive shift in the finite group
\[
 (C_2)^{\mathbb Z/n}\rtimes\mathbb Z/n.
\]
Conjugation of that lamp by the image of $f^\lambda$ shifts its
position by $-\lambda\pmod n$. If $f^\lambda$ centralizes $C$,
this forces $\lambda=0\pmod n$ for every $n$, hence
$\lambda=0$ in $\widehat{\mathbb Z}$. We have proved
\[
 C_{\widehat R}(C)\cap\widehat K=1.
\]
Conjugating by the entire actual gap word $a\in K$ gives
\begin{equation}\label{eq:B133-compact-cone-avoidance}
 C_{\widehat R}(c)\cap\widehat K=1.
\end{equation}
This uniform finite-wreath argument is the regional
cone-intersection proof; it uses every modulus, not a finite cutoff.

B.112 gives the full topology of the actual regional source $S$.
Its closure is consequently its embedded completion
$\widehat S\cong\widehat F_5$, and lies in
$C_{\widehat R}(c)$ by continuity. Equation
\eqref{eq:B133-compact-cone-avoidance} therefore gives
\[
 \widehat S\cap\widehat K=1.
\]
Conjugating the entire marked subgroup by $j$ gives
$j\widehat Sj\cap\widehat K=1$ in $\widehat R'$.
The actual base-winding character splits on $\langle k\rangle$,
so its closure is $\widehat{\mathbb Z}$, and its restriction to
that closure is an isomorphism. Since the character kills
$\widehat K$, it also proves
\begin{equation}\label{eq:B133-compact-source-avoidance}
 \overline{\langle k\rangle}\cap\widehat K=1,
 \qquad j\widehat Sj\cap\widehat K=1.
\end{equation}
These are intersections of actual closed subgroups. No completion
of a discrete intersection was substituted for them.

\emph{B.133.4. The applicable free profinite subgroup theorem.}
Lopes and Zalesskii's Corollary~2.18, with its preceding
Lemma~2.17, applies to two subgroups of the two vertices of a
proper free profinite amalgam when both subgroups meet the
amalgamated group trivially; their closed generated group is then
their free profinite product \cite[Corollary~2.18]{BLopesZalesskii2025}.
Apply it to \eqref{eq:B133-proper-completed-cover} and the two
closed subgroups in \eqref{eq:B133-compact-source-avoidance}.
It gives exactly
\[
 \overline{\mathscr F}^{\,\widehat D}
    =j\widehat Sj\amalg\overline{\langle k\rangle}
    \cong\widehat F_5\amalg\widehat{\mathbb Z}
    \cong\widehat F_6.
\]
The marked map from the completion of
\eqref{eq:B133-discrete-free-source} is this free profinite
product map. It is injective, with image the actual closure.
Thus $\mathscr F$ has its full induced topology in $D$ and in
$A$, proving \eqref{eq:B133-free-completion}.

The theorem is being used with trivial completed edge
intersections. It makes no assertion about the earlier carrier
with its nontrivial common subgroup generated by $P_5$, and no
general profinite Kurosh theorem is used.

\emph{B.133.5. The exact final source intersections.}
The discrete quotient in \eqref{eq:B133-pre-cross} has kernel
$N_\times$. Restricting it to the actual subgroup
$\mathscr F$ gives the first equality in
\eqref{eq:B133-kernel-intersections}. All six images have both
ordinary and mirror parity zero, so their equality in $X$ can
equivalently be tested in its actual subgroup $\Gamma_0$.

The universal property of finite quotients gives
\[
 \widehat X=\widehat A/\mathcal N_\times.
\]
The actual $\Gamma_0$ has finite index in $X$, and hence its
completion embeds in $\widehat X$ with its actual closed image.
The map from $\widehat{\mathscr F}$ lands in that image.
Using the established embedding of $\widehat{\mathscr F}$
in $\widehat A$, its kernel is therefore the literal
intersection $\widehat{\mathscr F}\cap\mathcal N_\times$.
This proves the second equality in
\eqref{eq:B133-kernel-intersections}.

One always has the inclusion
\[
 \overline{\mathscr F\cap N_\times}^{\,\widehat{\mathscr F}}
       \ \subseteq\ 
 \widehat{\mathscr F}\cap\mathcal N_\times.
\]
Equality is not inferred here. In particular this proof gives
neither full topology of the final source image nor an internal
normal-generation theorem for its kernel. The pre-cross full
topology and the post-cross noninjectivity are distinct conclusions.
\end{proof}

\begingroup\raggedright\footnotesize
The complete actual marking, topology, and source comparison is
retained in
\path{work_DE/next_marginal_source/pre_cross_source_full_topology27.md}.
The regional uniform cone-intersection proof is
\path{work_CE/annular_neck_cone_intersection23.txt}; the complete
pre-cross star and cofinal finite-kernel proof is
\path{work_A/ordinary_star_single_cross25.md}.
The publisher's complete Lemma~2.17 and Corollary~2.18, including
their proofs, were independently reopened at this application.
No new finite-group census or truncated source-orbit test is used.
\par\endgroup
